\documentclass[../../main.tex]{subfiles}

\begin{document}

\chapter[Top-Down Parsing]{Top-Down Parsing: Recursive Descent,\\ LL(1), FIRST, dan FOLLOW}
\label{chap:top-down}

\begin{outcomebox}
\textbf{Sub-CPMK:} \subcpmk{3.3}, \subcpmk{4.1}\\[3pt]
Setelah menyelesaikan bab ini, mahasiswa mampu:
\begin{enumerate}[label=\arabic*.,leftmargin=*]
  \item menjelaskan prinsip recursive descent dan predictive parsing;
  \item membuat tabel himpunan FIRST dan FOLLOW untuk grammar subset
        bahasa C;
  \item menguji keputusan parser LL(1) berdasarkan tabel tersebut;
  \item mengimplementasikan recursive descent parser dan predictive
        parser LL(1) dalam bahasa C;
  \item menangani kesalahan sintaksis dan menghasilkan keluaran
        validasi;
  \item menyelesaikan M2: parser untuk Subset-C.
\end{enumerate}
\end{outcomebox}

\section*{Materi Pokok}
\begin{itemize}[leftmargin=*]
  \item Recursive descent: satu fungsi per nonterminal.
  \item Himpunan FIRST dan FOLLOW serta tabelnya.
  \item Tabel parsing LL(1) dan keputusan LL(1).
  \item Implementasi recursive descent dan predictive parser dalam C.
  \item Penanganan kesalahan sintaksis, validasi, dan pemulihan.
\end{itemize}

\subfile{section-01}
\subfile{section-02}
\subfile{section-03}
\subfile{section-04}
\subfile{section-05}

\section*{Contoh}
\addcontentsline{toc}{section}{Contoh}

\begin{contoh}[Menghitung FIRST]
Untuk grammar
\[
E \to T E', \quad E' \to + T E' \mid \varepsilon, \quad
T \to F T', \quad T' \to * F T' \mid \varepsilon, \quad
F \to (E) \mid \code{id},
\]
FIRST(\(F\)) = \(\{\code{id}, (\}\). Karena \(T \to F T'\) dan
\(F\) tidak nullable, FIRST(\(T\)) = FIRST(\(F\)). Demikian pula
FIRST(\(E\)) = FIRST(\(T\)). Untuk \(E'\), FIRST(\(E'\)) =
\(\{+, \varepsilon\}\).
\end{contoh}

\begin{contoh}[Menghitung FOLLOW]
FOLLOW(\(E\)) = \(\{\$, )\}\) karena \(E\) adalah simbol awal dan
muncul di dalam kurung. FOLLOW(\(E'\)) sama dengan FOLLOW(\(E\)).
FOLLOW(\(T\)) = \(\{\$, ), +\}\) karena \(T\) diikuti \(E'\), dan
FIRST(\(E'\)) memuat \(+\). FOLLOW(\(F\)) = \(\{\$, ), *, +\}\)
karena \(F\) diikuti \(T'\), dan FIRST(\(T'\)) memuat \(*\).
\end{contoh}

\begin{contoh}[Mengisi tabel LL(1)]
Aturan \(E \to T E'\) diisikan pada kolom FIRST(\(T\)) =
\(\{\code{id}, (\}\). Aturan \(E' \to \varepsilon\) diisikan pada
kolom FOLLOW(\(E'\)) = \(\{\$, )\}\). Aturan
\(T' \to * F T'\) diisikan pada kolom \(*\). Sel yang tidak terisi
menandakan kesalahan sintaksis.
\end{contoh}

\begin{contoh}[Penelusuran stack]
Untuk masukan \(\code{id} + \code{id}\ \$\), stack dimulai dengan
\(\$E\). Aturan \(E \to TE'\) menghasilkan stack \(\$E'T\).
Aturan \(T \to FT'\) menghasilkan \(\$E'T'F\). Aturan
\(F \to \code{id}\) menghasilkan \(\$E'T'\code{id}\), lalu
\code{id} dicocokkan dan di-pop. Setelah \(T' \to \varepsilon\),
\(E' \to +TE'\), dan seterusnya, parser menerima masukan ketika
stack hanya berisi \(\$\).
\end{contoh}

\begin{contoh}[Pemulihan titik koma]
Untuk masukan
\begin{lstlisting}[language=C]
int a = 1
return 0;
\end{lstlisting}
parser melaporkan \code{baris 1, kolom 10: diharapkan ';' setelah
deklarasi, ditemukan 'return'}. Setelah itu, parser menyinkronkan
ke titik koma berikutnya dan melanjutkan. Program tetap dapat
dianalisis lebih lanjut meskipun mengandung kesalahan.
\end{contoh}

\section*{Aktivitas Pembelajaran dan Praktikum}
\addcontentsline{toc}{section}{Aktivitas Pembelajaran dan Praktikum}

\begin{table}[H]
\centering
\caption{Rencana aktivitas Bab X}
\label{tab:ch10-aktivitas}
\begin{tabularx}{\textwidth}{@{}llX@{}}
\toprule
\textbf{Sesi} & \textbf{Bentuk} & \textbf{Kegiatan} \\
\midrule
1 & Ceramah & Recursive descent dan predictive parsing \\
2 & Latihan & Menghitung FIRST dan FOLLOW \\
3 & Lokakarya & Membangun tabel LL(1) dan menelusuri stack \\
4 & Praktikum & M2: parser dan pengujian \\
\bottomrule
\end{tabularx}
\end{table}

\subsection*{M2 / Praktikum 2: Parser}
\textbf{Tujuan.} Membangun recursive descent parser untuk Subset-C
yang menghasilkan AST dan menangani kesalahan sintaksis.

\textbf{Keluaran.}
\begin{enumerate}[label=\arabic*.,leftmargin=*]
  \item Parser untuk ekspresi, deklarasi, dan pernyataan kontrol.
  \item AST minimal untuk ekspresi dan pernyataan.
  \item Minimal 20 kasus uji: valid dan tidak valid.
  \item Diagnostik sintaksis dengan lokasi dan harapan token.
  \item Laporan singkat.
\end{enumerate}

\textbf{Langkah kerja.}
\begin{enumerate}[label=\arabic*.,leftmargin=*]
  \item Hitung FIRST dan FOLLOW untuk grammar hasil Bab IX.
  \item Bangun tabel LL(1) dan periksa konflik.
  \item Implementasikan fungsi parser per nonterminal.
  \item Bangun AST saat mengurangi aturan.
  \item Tambahkan sinkronisasi dan pesan kesalahan.
  \item Uji dengan program valid dan tidak valid.
\end{enumerate}

\textbf{Rubrik.}
\begin{table}[H]
\centering
\caption{Rubrik Praktikum 2}
\label{tab:ch10-rubrik}
\begin{tabularx}{\textwidth}{@{}lXr@{}}
\toprule
\textbf{Aspek} & \textbf{Indikator} & \textbf{Bobot} \\
\midrule
Kebenaran & Parser menerima program valid, menolak yang tidak valid &
40\% \\
Kualitas kode & Struktur fungsi, penamaan, komentar & 30\% \\
Kelengkapan & Ekspresi, pernyataan, AST, diagnostik & 20\% \\
Pengujian & Kasus normal, batas, dan kesalahan & 10\% \\
\bottomrule
\end{tabularx}
\end{table}

\textbf{Definisi selesai.}
\begin{checklist}
  \item Parser lulus 20 kasus uji.
  \item AST terbentuk untuk program valid.
  \item Diagnostik menyebut lokasi dan harapan.
  \item Laporan diserahkan.
\end{checklist}

\section*{Latihan dan Refleksi}
\addcontentsline{toc}{section}{Latihan dan Refleksi}

\subsection*{Latihan}
\begin{enumerate}[label=\arabic*.,leftmargin=*]
  \item Hitung FIRST untuk setiap nonterminal pada grammar ekspresi.
  \item Hitung FOLLOW untuk setiap nonterminal pada grammar ekspresi.
  \item Bangun tabel LL(1) untuk grammar tersebut.
  \item Telusuri stack untuk masukan \code{id * id + id \$}.
  \item Jelaskan mengapa grammar dengan left recursion tidak dapat
        dipakai untuk recursive descent.
  \item Rancang fungsi \code{parse\_if} dan jelaskan pemulihannya.
  \item Berikan contoh kesalahan sintaksis dan pesan yang baik.
  \item Jelaskan perbedaan panic mode dan phrase-level recovery.
\end{enumerate}

\subsection*{Refleksi}
\begin{reflectionbox}
\begin{enumerate}[label=\arabic*.,leftmargin=*]
  \item Bagian mana dari perhitungan FIRST/FOLLOW yang paling sulit?
  \item Apakah Anda dapat menjelaskan mengapa parser LL(1) tidak
        pernah mundur?
  \item Bagaimana kualitas pesan kesalahan memengaruhi pengalaman
        pengguna?
\end{enumerate}
\end{reflectionbox}

\section*{Asesmen}
\addcontentsline{toc}{section}{Asesmen}

Bab X memuat Praktikum 2 (8\%) dan bagian UTS (20\%). Fokusnya adalah
FIRST/FOLLOW, LL(1), dan implementasi parser.

\begin{table}[H]
\centering
\caption{Instrumen asesmen Bab X}
\label{tab:ch10-asesmen}
\begin{tabularx}{\textwidth}{@{}lXl@{}}
\toprule
\textbf{Instrumen} & \textbf{Indikator} & \textbf{Bobot} \\
\midrule
Praktikum 2 & Parser, AST, 20 tes, diagnostik & 8\% \\
UTS & FIRST/FOLLOW, tabel LL(1), penelusuran & 20\% \\
Kuis & Recursive descent dan pemulihan & formatif \\
\bottomrule
\end{tabularx}
\end{table}

Kriteria kelulusan: parser menerima program valid, menolak program
tidak valid dengan pesan yang informatif, dan membangun AST yang
sesuai dengan struktur program.

\section*{Checklist Pencapaian Kompetensi}
\addcontentsline{toc}{section}{Checklist Pencapaian Kompetensi}

\begin{checklist}
  \item Saya dapat menjelaskan prinsip recursive descent.
  \item Saya dapat menghitung FIRST.
  \item Saya dapat menghitung FOLLOW.
  \item Saya dapat membuat tabel FIRST/FOLLOW.
  \item Saya dapat menentukan keputusan LL(1) dari tabel.
  \item Saya dapat membangun tabel LL(1).
  \item Saya dapat menelusuri predictive parser.
  \item Saya dapat mengimplementasikan recursive descent parser dan
        predictive parser dalam C.
  \item Saya dapat menangani kesalahan sintaksis.
  \item Saya dapat membangun AST minimal.
  \item Saya sudah menyelesaikan M2/Praktikum 2.
\end{checklist}

\section*{Rangkuman}
\addcontentsline{toc}{section}{Rangkuman}

Pada bab ini telah dibahas:
\begin{itemize}[leftmargin=*]
  \item recursive descent dan satu fungsi per nonterminal;
  \item himpunan FIRST dan FOLLOW serta cara menghitungnya;
  \item tabel parsing LL(1) dan predictive parser;
  \item konstruksi recursive descent yang sinkron dengan tabel;
  \item penanganan kesalahan sintaksis dan pemulihan.
\end{itemize}

Bab XI akan memperkenalkan parsing bottom-up, dimulai dari model
shift-reduce dan operator precedence. Pendekatan ini menangani
left recursion secara alami dan menjadi dasar parser LR pada Bab XII.

\section*{Evaluasi}
\addcontentsline{toc}{section}{Evaluasi}

\begin{enumerate}[label=\arabic*.,leftmargin=*]
  \item Jelaskan prinsip recursive descent.
  \item Hitung FIRST dan FOLLOW untuk grammar
        \(S \to A B\), \(A \to a \mid \varepsilon\),
        \(B \to b\).
  \item Bangun tabel LL(1) untuk grammar pada Soal 2.
  \item Telusuri stack untuk masukan \code{a b \$}.
  \item Mengapa grammar dengan left recursion harus ditransformasi?
  \item Jelaskan perbedaan panic mode dan error production.
  \item Rancang pesan kesalahan untuk deklarasi tanpa titik koma.
  \item Jelaskan hubungan antara tabel LL(1) dan kode recursive
        descent.
  \item Buat tabel FIRST/FOLLOW dan tentukan keputusan LL(1) untuk
        grammar yang diberikan; tunjukkan sel ganda bila ada.
\end{enumerate}

\subsection*{Kunci dan Rubrik Singkat}
\begin{itemize}[leftmargin=*]
  \item Soal 1--4: ketepatan perhitungan (70\%), penelusuran (30\%).
  \item Soal 5--6: pemahaman konsep (70\%), contoh (30\%).
  \item Soal 7--9: kualitas pesan, tabel, dan penjelasan (60\%),
        argumentasi (40\%).
\end{itemize}


\ifSubfilesClassLoaded{%
  \bibliographystyle{apalike}%
  \bibliography{../../references}%
}{}

\end{document}
